% ==============================================================
\section{結論}\label{sec:conclusion}
% ==============================================================

本論文では、Apriori-Window~\cite{aprioriwindow}を
任意の$d$次元グリッドに拡張する
\textbf{Multi-Dimensional Apriori-Window}を提案した。
固定空間始点ごとの1次元マルチセット還元、
重複考慮$t$軸ストライド、
および$(k{-}1)$-部分集合区間マップ交差枝刈りの3技法により、
全密集パターンとその任意形状の密集領域を厳密に列挙する。

実験により以下を示した：
(1)~4形状の合成データにおいて提案手法は厳密解を出力し、
既存手法はF1が最大0.917、Jaccardが最大0.190にとどまった；
(2)~$N = 10^4$--$10^8$で線形スケーリングを確認し、
$10^8$トランザクションを約9分で処理可能であった；
(3)~犯罪・苦情・地震の$d{=}3$--$4$実データにおいて、
複合犯罪ホットスポットや地震の深さ別密集など、
ドメイン知識と整合する解釈可能な示唆が得られた。

\paragraph{制約と限界.}
(1)~本手法はグリッド離散化を前提としており、
ビン幅および$\bm{W}$の選択が結果に影響する。
実際には、$\bm{W}$は対象とするパターンの時間的・空間的スケールに関するドメイン知識から初期化し、
感度分析で検証することができる；
$\theta$は推定されたバックグラウンド頻度を基準に設定することが実用的である。
(2)~$W_{\mathrm{sp}}$や$d$が増加すると
マルチセット構築の$O(N \cdot W_{\mathrm{sp}})$コストが増大し、
高次元設定では次元削減や粗いビン化が必要となる。
(3)~本手法は検出と記述に特化しており、
因果的説明や外部イベントとの関連付けは対象外である。

\paragraph{今後の課題.}
(1)~\emph{適応的パラメータ選択}：
局所的なデータ統計量（密度分布、空間的自己相関）に基づき
$\bm{W}$と$\theta$を自動調整するフレームワークは
実用上重要である。
(2)~\emph{更なる高速化}：
空間始点ごとの時間マルチセット構築が
主要な計算ボトルネックである。
累積和に基づく$O(1)$ウィンドウ集約や、
空間軸に沿った差分更新（スライス差分）による
冗長計算の排除が有望である。
